\chapter{External Data Products}
\label{chap:sp3-download}

\lupnt{} is not self-contained: almost every model in
Part~\ref{part:math} reads a numerical table that somebody else publishes and
maintains. Planetary positions come from the JPL Development Ephemerides
\citep{park2021de440} distributed as SPICE kernels \citep{acton1996spice};
Earth orientation and leap seconds come from the IERS conventions and the IERS
Earth Orientation Parameter (EOP) series \citep{iers2010}; spherical-harmonic
gravity fields come from the GRAIL and EGM analyses; and GNSS orbits and
clocks come from the IGS analysis centres, mirrored by NASA's Crustal Dynamics
Data Information System (CDDIS).

This chapter specifies the complete set of external products \lupnt{} expects,
where each one comes from, where it must land on disk, which environment
variable makes it findable, and which code path reads it. It then documents
the on-demand downloaders in detail --- the SPICE kernel refresh in
\file{cpp/lupnt/interfaces/spice.cc}, the IERS EOP refresh in
\file{cpp/lupnt/interfaces/eop.cc}, and the CDDIS GNSS product downloaders in
\file{cpp/lupnt/interfaces/sp3_loader.cc} and
\file{cpp/lupnt/interfaces/rinex_nav_loader.cc} --- because the GNSS ones need
credentials and are the most common source of first-run failures.

The numerical consequences of the choices described here (which DE version,
which EOP vintage, which lunar orientation kernel) are quantified in
\cref{chap:cross-validation}.

% ===========================================================================
\section{The data tree and the environment variables that find it}
\label{sec:sp3-data-tree}
% ===========================================================================

\lupnt{} resolves every data file through four functions in
\file{cpp/lupnt/core/file.cc}. They are the only places an environment
variable is read, so the contract is small and easy to check:

\begin{lstlisting}[language=cpp, caption={Path resolution.}, label={lst:sp3-file-cc}]
// cpp/lupnt/core/file.cc
std::filesystem::path GetBaseDir() {           // $LUPNT_PATH
  const char* base_path_env = std::getenv("LUPNT_PATH");
  LUPNT_CHECK(base_path_env, "Environment variable LUPNT_PATH is not set.", "File");
  return std::filesystem::path(base_path_env);
}

std::filesystem::path GetDataPath() { ... }    // $LUPNT_DATA_PATH, required

std::filesystem::path GetOutputDir(const std::string& dirname) {
  const char* output_path_env = std::getenv("LUPNT_OUTPUT_PATH");
  std::filesystem::path output_path;
  if (output_path_env == nullptr) output_path = GetDataPath() / "output";
  else                           output_path = std::filesystem::path(output_path_env);
  if (!dirname.empty()) output_path /= dirname;
  if (!std::filesystem::exists(output_path)) std::filesystem::create_directories(output_path);
  return output_path;
}

std::filesystem::path GetCspiceKernelDir() { return GetDataPath() / "ephemeris"; }
\end{lstlisting}

\cpp{GetFilePath(name)} performs a \emph{recursive} search of
\code{\$LUPNT\_DATA\_PATH} for a file whose name or stem matches \code{name}
(\cpp{FindFileInDir}), so most loaders name a file rather than a path: the
EOP loader asks for \code{eopc04\_08.62-now}, the IAU/SOFA loader for
\code{IAU\_SOFA.DAT}, and the gravity-field reader for
\code{grgm1200b.cof}. That search is convenient but it means \emph{two files
with the same basename anywhere under the data root are ambiguous}; keep the
tree as shipped.

\begin{table}[htbp]
  \centering
  \caption{Environment variables read by \lupnt{}. The \code{[activation]}
    block of \code{pixi.toml} sets \code{LUPNT\_DATA\_PATH},
    \code{LUPNT\_OUTPUT\_PATH}, \code{PECSIMPY\_BASE\_PATH} and
    \code{PYTHONPATH} automatically, so inside \code{pixi run} /
    \code{pixi shell} only the credentials need exporting.
    \code{LUPNT\_PATH} is not among them.}
  \label{tab:sp3-env}
  \small
  \begin{tabularx}{\textwidth}{@{}l W{1.35} W{0.65}@{}}
    \toprule
    Variable & Meaning & Default \\
    \midrule
    \code{LUPNT\_PATH}         & Repository root; \cpp{GetBaseDir()}. Required only by code that resolves repo-relative assets. & none (hard error) \\
    \code{LUPNT\_DATA\_PATH}   & Root of the reference-data tree, i.e.\ \code{data/LuPNT\_data}; \cpp{GetDataPath()}. Required by every loader. & none (hard error) \\
    \code{LUPNT\_OUTPUT\_PATH} & Root for generated products and download caches; \cpp{GetOutputDir()}. & \file{$LUPNT_DATA_PATH/output} \\
    \code{PECSIMPY\_BASE\_PATH}& Root of the plasma/GCPM coefficient tree, \code{data/LuPNT\_data/plasma}. & none \\
    \code{EARTHDATA\_USERNAME} & NASA Earthdata Login user, used by the CDDIS downloaders when \code{\textasciitilde/.netrc} is absent. & unset \\
    \code{EARTHDATA\_PASSWORD} & Matching password. & unset \\
    \code{LUPNT\_SKIP\_SPICE\_KERNEL\_DOWNLOAD} & Truthy value forces fully offline SPICE start-up: cached kernels only. & off \\
    \code{LUPNT\_SKIP\_SP3\_DOWNLOAD}  & Truthy value disables CDDIS SP3 fetches; only cache hits succeed. & off \\
    \code{LUPNT\_SKIP\_BRDC\_DOWNLOAD} & Same for RINEX broadcast navigation files. & off \\
    \code{LUPNT\_SKIP\_DEM\_DOWNLOAD}  & Same for the LOLA digital elevation models. & off \\
    \code{LUPNT\_TEST\_FIXTURES\_DIR}  & Overrides the C++ test fixture directory. & \code{cpp/test/fixtures} \\
    \bottomrule
  \end{tabularx}
\end{table}

The truthy test is shared by all the \code{LUPNT\_SKIP\_*} switches: any value
other than the empty string, \code{0}, \code{false}, \code{off} or \code{no}
(case-insensitive) disables the corresponding download.

\begin{lupntnote}
Outside \code{pixi}, run \file{scripts/config_development.sh} once. It appends
\code{LUPNT\_PATH}, \code{LUPNT\_DATA\_PATH} and \code{LUPNT\_OUTPUT\_PATH}
exports to the login shell profile, anchored at the repository it is run from.
\end{lupntnote}

\subsection{Layout on disk}
\label{sec:sp3-layout}

\begin{lstlisting}[language=shell, caption={The reference-data tree and the run-time cache.}, label={lst:sp3-tree}]
# $LUPNT_PATH is the repository root.
data/LuPNT_data/                 # $LUPNT_DATA_PATH  (gitignored, ~650 MB)
  ephemeris/                     #   SPICE kernels; GetCspiceKernelDir()
    ascii/                       #   ASCII DE440 header + coefficient files
  planetary_coeff/               #   EOP, leap seconds, IAU/SOFA CIP tables
  gravity/{earth,luna,mars,venus}/   # spherical-harmonic .cof files
  gnss/                          #   igs20.atx, gps_table.csv, galileo_table.xlsx
    sp3/                         #   scenario SP3 products (scripts/download_gnss_data.py)
    brdc/                        #   scenario RINEX navigation products
  antenna/                       #   transmit/receive gain patterns per constellation
  tle/<YYYY_MM_DD>/              #   dated TLE snapshots per constellation
  ground_station/                #   ground_stations.csv
  surface/                       #   lunar crater and PSR databases
  topo/                          #   body surface textures for plotting
  plasma/                        #   $PECSIMPY_BASE_PATH -- GCPM/plasmasphere tables
  dem/LOLA_5mpp/<site>/          #   LOLA DEM tiles, downloaded on demand
  examples/                      #   pre-computed example inputs

output/                          # $LUPNT_OUTPUT_PATH
  gnss_files/sp3/                #   GetOutputDir("gnss_files") -- SP3 download cache
  gnss_files/brdc/               #   RINEX navigation download cache
  ...                            #   one directory per application
\end{lstlisting}

The \file{data/LuPNT_data} tree is \emph{not} in version control. It is
published as a single \code{LuPNT\_data.zip} archive and pulled by whichever
front end runs first:

\begin{itemize}[leftmargin=1.4em]
  \item \file{cmake/FetchLuPNTData.cmake} downloads and extracts it during
    CMake configuration unless \code{LUPNT\_FETCH\_DATA} is \code{OFF}. It uses
    \file{data/LuPNT_data/gnss/igs20.atx} as the presence marker, so a partial
    tree missing that file triggers a re-download.
  \item \file{python/pylupnt/core/download_data.py} does the same from Python,
    keyed on \code{\$LUPNT\_DATA\_PATH/ephemeris} existing.
  \item \file{scripts/publish_lupnt_data.sh} pushes a new archive. The version
    string (\code{DATA\_VERSION} / \code{LUPNT\_DATA\_VERSION}) appears in
    three places and must be bumped together.
\end{itemize}

% ===========================================================================
\section{What \lupnt{} expects, and where it comes from}
\label{sec:sp3-products}
% ===========================================================================

\cref{tab:sp3-products} is the authoritative inventory. ``Refresh'' records
whether \lupnt{} tries to update the file from the network at run time
(\emph{auto}), whether a helper script fetches it on request (\emph{script}),
or whether it only ever arrives with the bundle (\emph{bundle}).

\begin{table}[htbp]
  \centering
  \caption{External data products, their provenance, and their location
    relative to \code{\$LUPNT\_DATA\_PATH} (or \code{\$LUPNT\_OUTPUT\_PATH}
    where marked).}
  \label{tab:sp3-products}
  \scriptsize
  \begin{tabularx}{\textwidth}{@{}L L l l@{}}
    \toprule
    Product (file) & Provenance & Location & Refresh \\
    \midrule
    \multicolumn{4}{@{}l}{\textit{SPICE kernels} \citep{acton1996spice}} \\
    \midrule
    Leap-second kernel \code{naif0012.tls} & NAIF generic kernels, \code{lsk/} & \code{ephemeris/} & bundle \\
    Planetary/lunar ephemeris \code{de440.bsp} & JPL DE440 \citep{park2021de440}, NAIF \code{spk/planets/} & \code{ephemeris/} & bundle \\
    Time ephemeris \code{de440t.bsp} & DE440t; adds NAIF body \code{1000000001}, the integrated $\TT-\TDB$ at the geocentre & \code{ephemeris/} & bundle \\
    Text PCK \code{pck00011.tpc} & NAIF planetary constants (radii, low-accuracy poles) & \code{ephemeris/} & bundle \\
    Earth orientation \code{earth\_latest\_high\_prec.bpc} & NAIF \code{pck/}; reconstructed $+$ short-term predicted ITRF93 & \code{ephemeris/} & auto \\
    Lunar orientation \code{moon\_pa\_de440\_200625.bpc} & NAIF \code{pck/}; DE440 principal-axes (PA) orientation & \code{ephemeris/} & auto \\
    Lunar frame kernel \code{moon\_de440\_250416.tf} & NAIF \code{fk/satellites/}; defines \code{MOON\_PA\_DE440} / \code{MOON\_ME\_DE440} & \code{ephemeris/} & auto \\
    \code{moon\_assoc\_pa.tf} & NAIF \code{fk/satellites/}; makes PA the default lunar body-fixed frame for legacy SPICE APIs & \code{ephemeris/} & auto \\
    Legacy fallbacks \code{earth\_000101\_241014\_240722.bpc}, \code{earth\_200101\_990825\_predict.bpc} & NAIF \code{pck/}, older snapshots & \code{ephemeris/} & bundle \\
    Older ephemerides \code{de405.bsp}, \code{mar097.bsp} & JPL/NAIF & \code{ephemeris/} & bundle \\
    ASCII DE440 (\code{header.440}, \code{ascp*.440}) & JPL ASCII distribution, read by \code{cpp/lupnt/interfaces/kernels.cc} & \code{ephemeris/ascii/de440/} & bundle \\
    \midrule
    \multicolumn{4}{@{}l}{\textit{Earth orientation, leap seconds, precession-nutation} \citep{iers2010}} \\
    \midrule
    Bundled EOP series \code{eopc04\_08.62-now} & IERS EOP 14 C04 (IAU1980) & \code{planetary\_coeff/} & bundle \\
    Latest EOP series \code{eopc04\_iers\_latest.txt} & \code{datacenter.iers.org}, \code{EOP\_14\_C04\_IAU1980\_one\_file\_1962-now.txt} & \code{planetary\_coeff/} & auto (on request) \\
    Leap-second table \code{tai-utc.dat} & IERS/USNO $\TAI-\UTC$ step table & \code{planetary\_coeff/} & bundle \\
    CIP/CIO table \code{IAU\_SOFA.DAT} & Tabulated IAU/SOFA $X$, $Y$, $s$ & \code{planetary\_coeff/} & bundle \\
    \midrule
    \multicolumn{4}{@{}l}{\textit{Gravity fields}} \\
    \midrule
    \code{grgm1200b.cof} (default lunar), \code{grgm900c.cof}, \code{LP165P.cof} & GRAIL-derived lunar fields \citep{lemoine2013grgm} & \code{gravity/luna/} & bundle \\
    \code{EGM96.cof} (default Earth), \code{JGM2.cof}, \code{JGM3.cof} & Earth geopotential models & \code{gravity/earth/} & bundle \\
    \code{GMM1.cof} (default Mars), \code{MGN75HSAAP.cof} (default Venus) & Mars/Venus geopotential models & \code{gravity/mars/}, \code{gravity/venus/} & bundle \\
    \midrule
    \multicolumn{4}{@{}l}{\textit{GNSS products}} \\
    \midrule
    SP3 precise orbits \code{COD0MGXFIN\_*\_01D\_05M\_ORB.SP3} & CDDIS \code{gnss/products/}; CODE MGEX final & \code{output/gnss\_files/sp3/} (cache) or \code{gnss/sp3/} (scenario) & script/auto \\
    RINEX~3 broadcast navigation \code{BRDC00IGS\_R\_*\_01D\_MN.rnx} & CDDIS \code{gnss/data/daily/}; IGS multi-GNSS merged & \code{output/gnss\_files/brdc/} or \code{gnss/brdc/} & script/auto \\
    ANTEX phase centres \code{igs20.atx} & IGS antenna exchange file, frame \code{IGS20} & \code{gnss/} & bundle \\
    Constellation tables \code{gps\_table.csv}, \code{galileo\_table.xlsx} & Per-satellite block/PRN metadata & \code{gnss/} & bundle \\
    TLE snapshots \code{gps\_<date>.txt}, \code{galileo\_<date>.txt}, \code{glonass\_<date>.txt}, \code{beidou\_<date>.txt}, \code{qzss\_<date>.txt}, \code{gnss\_<date>.txt} & Space-Track/CelesTrak two-line elements & \code{tle/2025\_01\_01/} and similar & bundle \\
    Antenna patterns & Per-constellation transmit and receive gain tables & \code{antenna/<name>/} & bundle \\
    \midrule
    \multicolumn{4}{@{}l}{\textit{Environment and surface}} \\
    \midrule
    LOLA 5\,m/px DEM tiles & \code{pgda.gsfc.nasa.gov/data/LOLA\_5mpp}; per-site GeoTIFF & \code{dem/LOLA\_5mpp/<site>/} & auto \\
    Kp index series & GFZ Potsdam JSON service & \code{plasma/} & auto \\
    GCPM/plasmasphere coefficients & Global Core Plasma Model tables & \code{plasma/data/} & bundle \\
    Ground-station catalogue \code{ground\_stations.csv} & Curated DSN/ESTRACK-style site list & \code{ground\_station/} & bundle \\
    Lunar crater / PSR databases & USGS crater database, LRO PSR list & \code{surface/} & bundle \\
    \bottomrule
  \end{tabularx}
\end{table}

\begin{lupntnote}
The Mars ephemeris is furnished under either of the two spellings in
circulation. NAIF ships the file --- and the \lupnt{} data bundle contains it
--- as \code{mar097.bsp}; the loader's existence check formerly tested only for
\code{mars097.bsp}, which never matched, so the kernel was silently never
furnished and Mars-centred states fell back on whatever \code{de440.bsp} alone
provides. The guard now iterates over both names and furnishes the first that
exists (\cref{sec:sp3-kernels}). Because the failure mode was a missing
\code{furnsh\_c} rather than an error, nothing in the log distinguished it from
a successful load, which is the general hazard of guarding a
\code{furnsh\_c} with a filename test.
\end{lupntnote}

% ===========================================================================
\section{SPICE kernels and the auto-refresh policy}
\label{sec:sp3-kernels}
% ===========================================================================

\cpp{spice::LoadSpiceKernel()} in \file{cpp/lupnt/interfaces/spice.cc} runs
once per process. It changes the working directory to
\cpp{GetCspiceKernelDir()} --- so every \cpp{furnsh\_c} argument below is a
bare filename resolved inside \file{data/LuPNT_data/ephemeris} --- and loads,
in order:

\begin{enumerate}[leftmargin=1.6em]
  \item \code{naif0012.tls}, the leap-second kernel. Without it no
    $\UTC$-based SPICE call works at all.
  \item \code{de440.bsp}, the planetary and lunar ephemeris
    \citep{park2021de440}. Its Chebyshev segments are additionally extracted
    into \lupnt{}'s own evaluator (\cpp{spk\_extract}, see
    \file{cpp/lupnt/interfaces/spice_cheby.cc}) so that body positions can be
    evaluated inside OpenMP regions without entering CSPICE, which is not
    thread safe.
  \item \code{de440t.bsp} \emph{if present}. This is DE440 plus one extra
    segment, NAIF body \code{1000000001}, holding the numerically integrated
    $\TT-\TDB$ at the geocentre. \lupnt{} prefers it over CSPICE's truncated
    analytic series in \cpp{unitim\_c}; the two differ by up to
    ${\sim}\SI{30}{\micro\second}$. If the file is absent a warning is emitted
    and the analytic series is used instead (see \cref{chap:time-conversions}).
  \item \code{pck00011.tpc}, the text planetary constants kernel. It only
    supplies the low-accuracy \code{IAU\_EARTH}/\code{IAU\_MOON} body-fixed
    frames, whose errors reach ${\sim}\SI{0.06}{\degree}$
    (${\sim}\SI{6}{\kilo\meter}$) for the Earth and ${\sim}\SI{0.005}{\degree}$
    (${\sim}\SI{155}{\meter}$) for the Moon.
  \item The high-accuracy orientation kernels, described next.
  \item The Mars ephemeris, \emph{if present}. The loader tries
    \code{mar097.bsp} and then \code{mars097.bsp}, furnishing the first that
    exists, and does nothing if neither does. Accepting both spellings is
    deliberate: NAIF publishes the file as \code{mar097.bsp}, but the longer
    form circulates widely enough that a guard on one name alone silently
    dropped the kernel.
\end{enumerate}

\begin{lstlisting}[language=cpp, caption={Furnishing the Mars ephemeris under either spelling.}, label={lst:sp3-mars}]
// cpp/lupnt/interfaces/spice.cc :: LoadSpiceKernel
for (const char* mars_bsp : {"mar097.bsp", "mars097.bsp"}) {
  if (std::filesystem::exists(mars_bsp)) {
    furnsh_c(mars_bsp);
    CheckSpiceFailure(fmt::format("Loading {}", mars_bsp));
    break;
  }
}
\end{lstlisting}

\subsection{Refreshing the orientation kernels}
\label{sec:sp3-orientation}

High-accuracy body-fixed frames --- \code{ITRF93} for the Earth and
\code{MOON\_PA\_DE440}/\code{MOON\_ME\_DE440} for the Moon --- require binary
PCKs that NAIF republishes as new data becomes available. \lupnt{} resolves
them at start-up rather than pinning versions:

\begin{itemize}[leftmargin=1.4em]
  \item \textbf{Earth.} NAIF publishes the current reconstructed plus
    short-term-predicted orientation under the stable name
    \code{earth\_latest\_high\_prec.bpc}, so \cpp{RefreshAndLoadKernel} simply
    re-downloads it every start-up. On failure the cached copy is used; if
    there is no cached copy, the older
    \code{earth\_200101\_990825\_predict.bpc} and
    \code{earth\_000101\_241014\_240722.bpc} are furnished instead.
  \item \textbf{Moon.} There is no stable name: each new ephemeris produces a
    new pair such as \code{moon\_pa\_de440\_200625.bpc} (binary PCK, NAIF
    \code{pck/}) and \code{moon\_de440\_250416.tf} (frame kernel, NAIF
    \code{fk/satellites/}). \lupnt{} fetches the two directory listings,
    matches the names against the regular expressions
    \code{moon\_pa\_de\textbackslash d+\_[0-9]\{6\}\textbackslash.bpc} and
    \code{moon\_de\textbackslash d+\_[0-9]\{6\}\textbackslash.tf}, and takes
    the lexicographically largest match --- the file names embed zero-padded DE
    numbers and dates, so lexicographic order is chronological order. If the
    listing cannot be retrieved, the newest locally cached match is used. The
    binary PCK must be loaded \emph{before} the frame kernel, since the FK is
    what binds the frame name to the orientation data.
  \item \code{moon\_assoc\_pa.tf} is then loaded so that the legacy
    (pre-N0062) SPICE APIs --- \code{LSPCN}, \code{ET2LST}, \code{ILLUM},
    \code{SRFXPT}, \code{SUBPT}, \code{SUBSOL} --- treat PA rather than ME as
    the lunar body-fixed frame.
\end{itemize}

Set \code{LUPNT\_SKIP\_SPICE\_KERNEL\_DOWNLOAD=1} to suppress all four network
lookups and run from the cache alone. This is the correct setting for
reproducible batch runs, because an auto-refreshed Earth kernel changes the
predicted portion of the orientation from week to week.

The distinction between the mean-Earth/polar-axis (ME) and principal-axes (PA)
lunar frames, and the rotation that connects them, is specified in
\cref{chap:frame-conversions}; \cref{chap:cross-validation} measures what the
choice is worth in metres.

% ===========================================================================
\section{Earth orientation parameters and leap seconds}
\label{sec:sp3-eop}
% ===========================================================================

The EOP table drives both the $\UTone-\UTC$ used by
\cref{chap:time-conversions} and the polar-motion and sidereal-rotation
matrices used by \cref{chap:frame-conversions}. It is loaded lazily by
\cpp{GetEopFileData()} / \cpp{GetEopData(mjd\_utc)} in
\file{cpp/lupnt/interfaces/eop.cc}, from
\cpp{GetFilePath(EOP\_FILENAME)} where \code{EOP\_FILENAME} is
\code{"eopc04\_08.62-now"} (declared in \file{cpp/lupnt/core/constants.h}).

The file is an IERS EOP 14 C04 (IAU1980) series: fourteen header lines
followed by whitespace-separated columns holding the year, month, day, MJD in
$\UTC$, the pole coordinates $x_p$ and $y_p$, $\UTone-\UTC$, the length of day,
the nutation corrections $\delta\psi$ and $\delta\epsilon$, and the formal
error of each. Values are interpolated with a third-order Lagrange rule;
outside the tabulated span the nearest endpoint is held constant rather than
extrapolated.

\begin{lstlisting}[language=cpp, caption={On-demand refresh from the IERS data centre.}, label={lst:sp3-eop}]
// cpp/lupnt/interfaces/eop.cc :: LoadLatestEopFromIers
static const std::string url
    = "https://datacenter.iers.org/data/latestVersion/"
      "EOP_14_C04_IAU1980_one_file_1962-now.txt";
std::filesystem::path dest_path = GetDataPath() / "planetary_coeff" / "eopc04_iers_latest.txt";
if (DownloadFile(url, dest_path)) { LoadEopFileData(dest_path, force); return true; }
// otherwise: warn and fall back to the bundled snapshot
\end{lstlisting}

\begin{lupntwarning}
\cpp{LoadLatestEopFromIers} replaces the global EOP table for the whole
process. The GMAT and Orekit cross-validation tolerances in
\cref{chap:cross-validation} are calibrated against the \emph{bundled}
vintage, which is why \code{data.eop\_latest\_iers} restores
\cpp{GetFilePath(EOP\_FILENAME)} before it returns. Anything that refreshes
EOP mid-run and then compares against a fixture will see tens of milliseconds
of $\UTone$ shift, i.e.\ hundreds of metres at geostationary radius.
\end{lupntwarning}

The separate \code{tai-utc.dat} table gives the integer $\TAI-\UTC$ steps and
\code{IAU\_SOFA.DAT} the tabulated CIP coordinates $X$, $Y$ and the CIO locator
$s$; both are read through the same \cpp{GetFilePath} name lookup.

% ===========================================================================
\section{GNSS products: SP3, clocks, RINEX navigation, ANTEX}
\label{sec:sp3-gnss}
% ===========================================================================

\subsection{What the products are}
\label{sec:sp3-formats}

\paragraph{SP3 (Standard Product 3).} A plain-text tabulation of GNSS
satellite positions --- and, in the \code{c}/\code{d} variants, clock offsets
--- on a uniform epoch grid. Each epoch begins with a line
\code{* yyyy mm dd hh mm ss.sssssss} in GPS time, followed by one
\code{P<sv>} record per satellite giving $x$, $y$, $z$ in kilometres in the
terrestrial (ECEF) frame and the clock offset in microseconds, in fixed
columns. The value \code{999999.999999} is the SP3 sentinel for
``clock unavailable''. \lupnt{}'s parser converts positions to metres,
converts the epoch from GPS to $\TAI$, and maps any clock above
\num{999990} to NaN:

\begin{lstlisting}[language=cpp, caption={SP3 record parsing.}, label={lst:sp3-parse}]
// cpp/lupnt/interfaces/sp3_loader.cc :: Sp3Loader::ParseFile
current_epoch = ConvertTime(GregorianToTime(year, month, day, hour, minute, second),
                            Time::GPS, Time::TAI);
double x = std::stod(line.substr(4, 14)) * 1000.0;   // km -> m
double y = std::stod(line.substr(18, 14)) * 1000.0;  // km -> m
double z = std::stod(line.substr(32, 14)) * 1000.0;  // km -> m
double clock_us = std::stod(line.substr(46, 14));
double clock = (clock_us > kSp3ClockSentinelUs) ? std::numeric_limits<double>::quiet_NaN()
                                                : clock_us * 1e-6;
\end{lstlisting}

\paragraph{Clocks.} \lupnt{} does \emph{not} consume separate RINEX clock
(\code{.CLK}) products. Satellite clock offsets are taken from the fourth
column of the SP3 position records, at the SP3 sampling interval
(\SI{300}{\second} for the CODE MGEX final product). Sub-\SI{300}{\second}
clock structure --- which a \SI{30}{\second} \code{.CLK} product would
resolve --- is therefore not modelled.

\paragraph{RINEX navigation (broadcast ephemeris).} The Keplerian-plus-
correction parameter sets broadcast by the satellites themselves, in RINEX~3
mixed-navigation form (\code{MN}). These are what a real receiver would use;
comparing them against SP3 is the standard way to visualise broadcast orbit
and clock error, and \file{python/examples/ex3_gnss_interface.ipynb} does
exactly that.

\paragraph{ANTEX.} Antenna phase-centre offsets and variations, per satellite
block and per frequency, in the IGS \code{ANTEX} format. \lupnt{} ships
\file{data/LuPNT_data/gnss/igs20.atx} (frame \code{IGS20}) and reads it with
\cpp{AntexLoader}; a phase-centre offset is what converts the SP3
centre-of-mass position into the antenna phase-centre position that the
observation equation actually needs.

\subsection{Provider and naming conventions}
\label{sec:sp3-naming}

Both downloaders derive a canonical product name from the query epoch. The
epoch is first converted to GPS time, the GPS week and day-of-week are
computed from the GPS epoch 1980-01-06, and the calendar year and day-of-year
are taken from the corresponding Modified Julian Date:

\begin{lstlisting}[language=cpp, caption={Product naming.}, label={lst:sp3-naming}]
// cpp/lupnt/interfaces/sp3_loader.cc :: BuildSp3ProductInfo
info.gps_week   = (int)std::floor(delta / kSecsWeek);              // delta = t_gps - t_gps_epoch
info.day_of_week= (int)std::floor((delta - info.gps_week * kSecsWeek) / SECS_DAY);
LUPNT_CHECK(info.gps_week >= 1962, "... modern CDDIS COD/MGEX `.SP3.gz` products only", "Sp3Loader");
info.filename = fmt::format("COD0MGXFIN_{:04d}{:03d}0000_01D_05M_ORB.SP3", info.year, info.doy);
info.zipped_filename = info.filename + ".gz";
info.url = fmt::format("https://cddis.nasa.gov/archive/gnss/products/{:04d}/{}",
                       info.gps_week, info.zipped_filename);
\end{lstlisting}

\begin{table}[htbp]
  \centering
  \caption{CDDIS products \lupnt{} downloads, with the naming template and the
    archive path. \code{YYYY} is the calendar year, \code{DDD} the
    zero-padded day of year, \code{YY} the two-digit year, and \code{WWWW}
    the zero-padded GPS week.}
  \label{tab:sp3-naming}
  \small
  \begin{tabularx}{\textwidth}{@{}l L L@{}}
    \toprule
    Product & File name template & CDDIS URL \\
    \midrule
    Precise orbits (SP3) &
      \code{COD0MGXFIN\_YYYYDDD0000\_01D\_05M\_ORB.SP3.gz} &
      \code{https://cddis.nasa.gov/archive/gnss/products/WWWW/} \\
    Broadcast navigation (RINEX~3), primary &
      \code{BRDC00IGS\_R\_YYYYDDD0000\_01D\_MN.rnx.gz} &
      \code{.../archive/gnss/data/daily/YYYY/DDD/YYp/} \\
    Broadcast navigation, fallback &
      same &
      \code{.../archive/gnss/data/daily/YYYY/brdc/} \\
    \bottomrule
  \end{tabularx}
\end{table}

Decoding the SP3 template: \code{COD} is the CODE analysis centre (Center for
Orbit Determination in Europe, Bern), \code{0} the solution number,
\code{MGX} the multi-GNSS experiment campaign (GPS, GLONASS, Galileo, BeiDou
and QZSS in one file), \code{FIN} the final --- as opposed to rapid or
ultra-rapid --- product, \code{01D} the one-day file span, \code{05M} the
\SI{300}{\second} sampling, and \code{ORB} the orbit content. The
\code{BRDC00IGS\_R} template is the IGS merged multi-GNSS broadcast file,
\code{R} marking a data-centre-generated (rather than station) product and
\code{MN} mixed navigation.

Two constraints follow from the naming scheme. Epochs before the GPS epoch are
rejected outright, and epochs before GPS week 1962 are rejected because the
long-form CDDIS naming convention (and the MGEX final product) does not exist
that far back. Products are daily, so a request is quantised to a calendar
day; helper scripts anchor at 12:00 to keep a midnight rollover from selecting
the neighbouring file.

\begin{lupntnote}
Both loaders accept several files at once. \cpp{Sp3Loader::LoadFile} sorts and
de-duplicates the merged epoch table per satellite: at a day boundary the same
epoch appears in two files, typically with the sentinel clock in the earlier
file and the real value in the later one, so the record with a valid clock
wins and, failing that, the later file wins. A run that crosses midnight
should load both days.
\end{lupntnote}

% ===========================================================================
\section{NASA Earthdata credentials}
\label{sec:sp3-earthdata}
% ===========================================================================

CDDIS is behind NASA's Earthdata Login (URS). Before the first download:

\begin{enumerate}[leftmargin=1.6em]
  \item Create an account at \url{https://urs.earthdata.nasa.gov}.
  \item Review the Earthdata Login API and integration notes at
    \url{https://www.earthdata.nasa.gov/engage/open-data-services-software/earthdata-developer-portal/earthdata-login-api}.
  \item Store the credentials in \file{~/.netrc} --- the file is
    \code{.netrc} in the \emph{home} directory, not \code{./netrc} in the
    repository:
\end{enumerate}

\begin{lstlisting}[language=shell, caption={Credential setup.}, label={lst:sp3-netrc}]
touch ~/.netrc
chmod 0600 ~/.netrc
cat >> ~/.netrc <<'EOF'
machine urs.earthdata.nasa.gov
  login YOUR_EARTHDATA_USERNAME
  password YOUR_EARTHDATA_PASSWORD
EOF
\end{lstlisting}

The \code{0600} mode matters: many HTTP clients refuse a world-readable
\code{.netrc}. If this is the first scripted CDDIS access, sign in through a
browser once and approve the CDDIS application-access prompt --- an
unapproved application yields a login page rather than a file, which is the
failure mode described in \cref{sec:sp3-troubleshooting}.

For short tests the loaders also accept credentials from the environment:

\begin{lstlisting}[language=shell]
export EARTHDATA_USERNAME="YOUR_EARTHDATA_USERNAME"
export EARTHDATA_PASSWORD="YOUR_EARTHDATA_PASSWORD"
\end{lstlisting}

Prefer \file{~/.netrc} for routine use, so the password is not exposed in
shell history, process environments, or terminal logs.

\subsection{How the transfer is actually performed}
\label{sec:sp3-curl}

The downloaders shell out to \code{curl} rather than linking an HTTP client:

\begin{lstlisting}[language=cpp, caption={The CDDIS fetch.}, label={lst:sp3-curl}]
// cpp/lupnt/interfaces/sp3_loader.cc :: DownloadToFileEarthdata
const bool has_env_auth = std::getenv("EARTHDATA_USERNAME") != nullptr
                          && std::getenv("EARTHDATA_PASSWORD") != nullptr;
std::string auth_arg;
if (has_env_auth) auth_arg = "-u \"$EARTHDATA_USERNAME:$EARTHDATA_PASSWORD\" ";
// The URS OAuth login redirects cddis -> urs.earthdata -> cddis/proxyauth -> original URL and
// hands off session state via a cookie.  Without -c/-b persisting that cookie across the
// redirect chain the final hop restarts the OAuth dance and loops until curl's 50-redirect
// limit, even with correct credentials.
std::string cmd = fmt::format(
    "curl -fsSL --netrc-optional -c {0} -b {0} {1}--connect-timeout 10 --max-time 600 -o {2} {3}",
    ShellQuote(cookie_jar), auth_arg, ShellQuote(tmp_path.string()), ShellQuote(url));
\end{lstlisting}

The cookie jar is \file{~/.urs_cookies} when \code{HOME} is set, otherwise
\code{.urs\_cookies} beside the destination. Downloads land in a
\code{.part} temporary and are renamed only after a non-empty transfer, so an
interrupted download never leaves a truncated product in the cache. The
gzipped product is then expanded with \code{gzip -dc} and the archive removed.
Both \code{curl} and \code{gzip} must be on \code{PATH}.

Before decompressing, \cpp{LooksLikeHtml} inspects the first 256 bytes: if
they begin with \code{<!doctype html} or \code{<html},
the download is deleted and an explicit error is raised, because CDDIS
returned the Earthdata Login page instead of the file.

% ===========================================================================
\section{Download scripts and pixi tasks}
\label{sec:sp3-scripts}
% ===========================================================================

\begin{table}[htbp]
  \centering
  \caption{The download entry points that exist in the repository, and where
    each one writes.}
  \label{tab:sp3-scripts}
  \small
  \begin{tabularx}{\textwidth}{@{}L L L@{}}
    \toprule
    Entry point & Purpose & Destination \\
    \midrule
    \code{scripts/download\_gnss\_data.py} &
      Scenario data for \code{configs/lunar\_gnss\_odts.yaml} and the
      \code{ex6\_gnss\_odts} examples. Defaults to the epoch in the YAML and
      fetches \code{--days} consecutive daily products (default 2). &
      \code{data/LuPNT\_data/gnss/sp3}, \code{data/LuPNT\_data/gnss/brdc} \\
    \code{scripts/download\_gnss\_test\_fixtures.py} \newline
      (\code{pixi run download-gnss-test-data}) &
      The fixtures the C++ suite parses: SP3 for 2026-01-14 and 2026-01-15,
      BRDC for 2026-01-14. Uses the loaders' \emph{default} cache directory. &
      \code{output/gnss\_files/sp3}, \code{output/gnss\_files/brdc} \\
    \code{python/pylupnt/core/download\_data.py} &
      Fetches and unpacks the whole \code{LuPNT\_data.zip} bundle. &
      \code{data/LuPNT\_data} \\
    \code{cmake/FetchLuPNTData.cmake} &
      Same bundle, during CMake configuration
      (\code{-DLUPNT\_FETCH\_DATA=OFF} to disable). &
      \code{data/LuPNT\_data} \\
    \code{scripts/publish\_lupnt\_data.sh} &
      Publishes a new bundle version (maintainers). & --- \\
    \bottomrule
  \end{tabularx}
\end{table}

\begin{lstlisting}[language=shell, caption={Typical first-run sequence.}, label={lst:sp3-run}]
# 1. Reference-data bundle (also happens automatically at CMake configure time).
pixi run python -c "import pylupnt"          # triggers download_data.py if needed

# 2. GNSS products for the shipped scenario (needs Earthdata credentials).
pixi run python scripts/download_gnss_data.py
pixi run python scripts/download_gnss_data.py --epoch 2026-01-01T02:00:00 --days 2

# 3. GNSS fixtures for the C++ test suite.
pixi run download-gnss-test-data

# 4. Examples that download on demand.
pixi run build-examples                      # builds cpp/examples/tutorials/ex3, ex5
\end{lstlisting}

The two GNSS destinations are deliberately different.
\file{output/gnss_files} is the loaders' own cache --- it is what
\cpp{GetOutputDir("gnss\_files")} returns and what the \cpp{GnssFilesDir()}
helper in \file{cpp/test/interfaces/test_sp3_loader.cc} resolves to --- while
\file{data/LuPNT_data/gnss/sp3} and \file{data/LuPNT_data/gnss/brdc} are the
scenario input directories referenced from the YAML configuration. A file fetched into either is reusable
by both the C++ and Python interfaces, because the Python bindings are the C++
loaders.

The tests that need those fixtures are excluded from the CI subset
(\code{pixi run test-cpp-ci}), so a machine without Earthdata credentials
still gets a green run; \code{pixi run test-cpp} runs them.

\begin{lupntwarning}
The Sphinx page this chapter derives from refers to
\code{pixi run download-sp3-example} and
\code{pixi run download-sp3-example-cpp}. No such tasks exist in
\file{pixi.toml}. The equivalents that do exist are
\code{pixi run download-gnss-test-data} and the notebooks
\file{python/examples/ex3_gnss_interface.ipynb} and
\file{python/examples/ex5_gnss_measurement_sim.ipynb}, whose C++ counterparts
are \file{cpp/examples/tutorials/ex3_gnss_interface.cc} and
\file{cpp/examples/tutorials/ex5_gnss_measurement_sim.cc}.
\end{lupntwarning}

% ===========================================================================
\section{Using the loaders}
\label{sec:sp3-api}
% ===========================================================================

The download helpers are static: they compute the product name, fetch and
cache the file if it is missing, and return the local path. Parsing is then a
constructor call. \cpp{Sp3Loader} additionally exposes
\cpp{FilenameForEpoch} and \cpp{UrlForEpoch}, which do the naming arithmetic
without touching the network --- useful for dry runs and for building an
offline mirror.

\begin{lstlisting}[language=py, caption={Python interface.}, label={lst:sp3-python}]
# python/examples/ex3_gnss_interface.ipynb (condensed)
import pylupnt as pnt

# A TAI epoch whose daily product we want.
t_tai = pnt.convert_time(
    pnt.gregorian_to_time(2026, 1, 1, 12, 0, 0), pnt.Time.TDB, pnt.Time.TAI
)

# Download + cache the CODE MGEX final SP3 (Earthdata Login required), then parse it.
sp3_path = pnt.Sp3Loader.download_file_for_epoch(t_tai, pnt.Time.TAI)
sp3 = pnt.Sp3Loader(sp3_path)

print(sp3.get_satellites()[:8])          # ['C06', 'C07', ..., 'E02', ...]
print(sp3.get_time_span("G01"))          # (t_first_tai, t_last_tai)
rv_ecef, clock_bias_s = sp3.get_pos_vel_clock("G01", t_tai)   # [m, m/s], [s]

# Broadcast ephemeris and antenna phase centres use the same pattern.
brdc = pnt.RinexNavLoader(pnt.RinexNavLoader.download_file_for_epoch(t_tai, pnt.Time.TAI))
antex = pnt.AntexLoader(pnt.get_file_path("igs20.atx"))
\end{lstlisting}

The optional third argument of \py{download\_file\_for\_epoch} is a cache
directory; passing one is how \file{scripts/download_gnss_data.py} redirects
the fetch into \file{data/LuPNT_data/gnss}.

\begin{lstlisting}[language=cpp, caption={C++ interface.}, label={lst:sp3-cpp}]
// cpp/examples/tutorials/ex5_gnss_measurement_sim.cc (condensed)
#include <lupnt/lupnt.h>
using namespace lupnt;

Real base_day_tai = ConvertTime(GregorianToTime(2026, 1, 14, 12, 0, 0.0), Time::TDB, Time::TAI);

std::vector<std::filesystem::path> sp3_paths;
for (int d = 0; d < 2; ++d)
  sp3_paths.push_back(Sp3Loader::DownloadFileForEpoch(base_day_tai + d * 86400.0, Time::TAI));

Sp3Loader loader(sp3_paths);          // merges, sorts and de-duplicates both days
const auto& sats = loader.GetSatellites();
\end{lstlisting}

\begin{lupntnote}
\cpp{Sp3Loader::GetPosVelClock} does not simply interpolate the tabulated
samples. Positions are fitted with per-satellite Chebyshev segments (segment
length $8\times$ the minimum sample step, up to 12 coefficients), and the
segments are sampled from the tabulated arc with a local high-order Lagrange
interpolant rather than a linear one --- over a \SI{300}{\second} product
interval a GNSS arc is curved enough that linear sampling would inject
kilometre-level chord error that the Chebyshev fit would then reproduce
faithfully. Clocks are interpolated linearly at query time, deliberately, to
avoid Gibbs-like ringing at day boundaries.
\end{lupntnote}

% ===========================================================================
\section{Troubleshooting}
\label{sec:sp3-troubleshooting}
% ===========================================================================

\begin{table}[htbp]
  \centering
  \caption{Common failures and their causes.}
  \label{tab:sp3-troubleshooting}
  \small
  \begin{tabularx}{\textwidth}{@{}L L@{}}
    \toprule
    Symptom & Cause and remedy \\
    \midrule
    ``CDDIS returned an Earthdata Login HTML page instead of the requested SP3 file'' &
      The request reached URS but was not authenticated. Check that
      \code{\textasciitilde/.netrc} exists in the home directory (not the repository), that
      the machine entry is exactly \code{urs.earthdata.nasa.gov}, that the
      mode is \code{0600}, that the credentials are correct, and that CDDIS
      application access has been approved in a browser. \\
    ``Maximum (50) redirects followed'' &
      The URS cookie is not being persisted. \lupnt{}'s own \code{curl}
      invocation passes \code{-c}/\code{-b}; a hand-written command must do
      the same. \\
    ``Failed to download SP3 file from CDDIS'' with no HTML &
      Network, proxy, or timeout (\code{--connect-timeout 10},
      \code{--max-time 600}), or \code{LUPNT\_SKIP\_SP3\_DOWNLOAD} is set. \\
    ``C++ SP3 download helper supports modern CDDIS COD/MGEX \code{.SP3.gz} products (GPS week $\geq$ 1962) only'' &
      The requested epoch predates the long-form naming convention. Supply the
      file manually and construct \cpp{Sp3Loader} from the path. \\
    ``Environment variable LUPNT\_DATA\_PATH is not set'' &
      Running a binary outside \code{pixi}. Export the variables in
      \cref{tab:sp3-env} or run \code{scripts/config\_development.sh}. \\
    \code{File not found: <name>} from \cpp{GetFilePath} &
      The bundle is missing or partial. Delete \code{data/LuPNT\_data} and let
      CMake or \code{download\_data.py} re-fetch it. \\
    SPICE warns that \code{de440t.bsp} was not found &
      $\TT \leftrightarrow \TDB$ silently degrades to the analytic series;
      see \cref{chap:time-conversions}. \\
    \bottomrule
  \end{tabularx}
\end{table}

The credential path can be tested independently of \lupnt{}:

\begin{lstlisting}[language=shell]
# -L follows the Earthdata redirect chain, -n reads ~/.netrc,
# -c/-b persist the URS session cookie across the chain.
curl -L -n -c ~/.urs_cookies -b ~/.urs_cookies -O \
  "https://cddis.nasa.gov/archive/gnss/products/2399/COD0MGXFIN_20260010000_01D_05M_ORB.SP3.gz"
\end{lstlisting}

% ===========================================================================
\section{Model boundaries}
\label{sec:sp3-boundaries}
% ===========================================================================

\begin{itemize}[leftmargin=1.4em]
  \item \textbf{One SP3 product only.} The downloader hard-codes the CODE
    MGEX \emph{final} product. Rapid and ultra-rapid products, other analysis
    centres, and the legacy short file names are not supported by
    \cpp{BuildSp3ProductInfo}; such files can still be parsed by passing an
    explicit path.
  \item \textbf{No separate clock products.} Clocks come from the SP3 records
    at \SI{300}{\second}. Applications sensitive to shorter-term clock
    behaviour need an external \code{.CLK} source.
  \item \textbf{No SP3 velocity records.} Only \code{P} records are parsed;
    velocities are obtained by differentiating the Chebyshev position fit, not
    from \code{V} records.
  \item \textbf{Ephemeris version is fixed by the bundle.} \lupnt{} evaluates
    DE440 \citep{park2021de440}. A comparison against a tool that evaluated a
    different DE version would carry the full inter-ephemeris difference, so
    the cross-validation avoids it by pointing GMAT at \lupnt{}'s own
    \file{de440.bsp}: with every tool on DE440 no DE-version offset remains,
    as quantified in \cref{chap:cross-validation}.
  \item \textbf{Orientation kernels are time-varying inputs.} The Earth binary
    PCK is re-downloaded at every start-up and its predicted tail changes
    between releases. Pin the data by setting
    \code{LUPNT\_SKIP\_SPICE\_KERNEL\_DOWNLOAD=1} when bit-reproducibility
    matters.
  \item \textbf{EOP outside the table is held, not extrapolated.} Epochs
    beyond the last tabulated day reuse the endpoint values, so polar motion
    and $\UTone-\UTC$ freeze rather than diverge --- convenient, but the error
    grows without any diagnostic. Refresh the series for far-future epochs.
  \item \textbf{Linear interpolation across leap seconds.} The daily
    $\UTone-\UTC$ column is interpolated straight through the \SI{1}{\second}
    step, which corrupts $\UTone$ within roughly a day of an insertion. The
    measured size of the effect, and the way the cross-validation tests handle
    it, are given in \cref{chap:cross-validation}.
  \item \textbf{External tooling assumed.} \code{curl} and \code{gzip} must be
    installed; there is no in-process HTTP or inflate implementation.
\end{itemize}
